\section{Mechanisation}\label{sec:mechanisation}

The development is about 6{,}000 lines of Rocq 9.0.1, available at
\url{\repobase}, using Mathematical
Components~\cite{mathcomp} and, for the coinductive definitions of the original
paper, paco~\cite{paco} (\cref{tab:files}).

\begin{table}[t]
\centering
\caption{The development.}
\label{tab:files}
\footnotesize
\begin{tabularx}{\columnwidth}{@{}lrL@{}}
\toprule
file & lines & contents \\
\midrule
\mechfile{theories/definitions.v}{definitions.v} & 908 & the calculus, traces, $\NI$, the equivalences \\
\mechfile{theories/theorems.v}{theorems.v} & 1337 & the composition theorems \\
\mechfile{theories/models.v}{models.v} & 1079 & the generic model, both policies, \cref{thm:noop}, example traces \\
\mechfile{theories/noninterference.v}{noninterference.v} & 1749 & the model's equivalences, \cref{thm:leak,thm:main} \\
\mechfile{theories/adequacy.v}{adequacy.v} & 904 & the original definitions, the equivalences with ours, \cref{thm:leak,thm:main} in the original's terms \\
\bottomrule
\end{tabularx}
\end{table}

\paragraph{Types} Input and output types range over an inductive universe \m{Ty} of
codes, interpreted into \code{Set}, rather than over \code{Set} directly. This makes
reductions invertible, which the negative result needs: refuting non-interference
means showing that a trace is \emph{not} admitted.

\paragraph{Axioms} The development assumes nothing beyond the two parameters
that fix the security lattice: no classical logic, and no uniqueness of identity
proofs. The latter is what dependent elimination on the type-indexed reduction
relations would need; we invert reductions through determinism of reduction
instead, one lemma per rule (\m{inv_mapO} and its siblings).

\subsection{Non-interference over finite traces}\label{sec:departures}

Our definitions follow Rafnsson et al., with one essential departure. The original
defines non-interference as a coinductive simulation over infinite streams of
labels, whose clauses apply at the head of a stream. We quantify over finite traces
instead, with closure clauses that apply at an arbitrary position (Definition 1).
This makes every proof in the development an induction on a trace.

The original definition is written out in \mechfile{theories/adequacy.v}{adequacy.v},
and \m{NI_paper} proves the two equivalent. The proof passes through infinite
streams (\m{NI_SNI}) and rests on two properties: finite and infinite behaviour
determine each other, which needs reduction to be deterministic (to assemble
finite runs into one infinite run) and never to block (to continue a finite run);
and insertion commutes with taking prefixes. Obliviousness, the other definition
that quantifies over behaviour, is likewise stated over traces and equivalent to
the original (\m{ObliviousDis_iff}). As a consequence, \cref{thm:leak,thm:main}
also hold with non-interference read as in the original, hypotheses included
(\m{model_immediate_not_NI_paper}, \m{model_sliced_NI_paper}).

One further difference is only notational. The original adjoins the distinguished
value $\unobs$ to every type; we keep the values as they are and make
unobservability a predicate, tied to the equivalence by the characterisation law.
\m{NI_lEquiv} shows that the two give the same notion of non-interference, and
this paper uses the original notation throughout.

The composition theorems themselves are proved directly for our definitions.
